% Ad Familiares 15.6 --- headnote
% ad-familiares-15-06, late July 50 BC (shortly before 1 August), Tarsus

\textit{Cicero from Tarsus to Cato at Rome, written
shortly before the Kalends of Sextilis 50 BC (the
manuscript dateline: \textit{Scr. Tarsi paulo ante in.
K. Sext. a. 704 (50)}). The reply to Cato's letter
(Fam. 15.5), which conceded a citation for integrity but
withheld support for the public thanksgiving. Cicero had
asked Cato for both in Fam. 15.4; the Senate had in the
event voted the \textit{supplicatio} despite Cato's
dissent. By the time Cicero composes this answer he
knows that he has carried the day on the substance; what
remains is the management of the friendship.}

\textit{The architecture of the reply mirrors Cato's,
and slightly twists it. Section 1 chooses the high road,
opening with a verse from Naevius --- Hector to his
father, ``I am glad to be praised by you, a praised man''
--- and praising Cato's testimony as worth more than any
chariot or laurel: the highest tribute in the corpus,
delivered to the one senator who had argued against him.
Section 2 turns Cato's conditional construction back
upon him with deliberate care. Cicero declines Cato's
word ``craving'' for his ``wish''; concedes that an
honour ought not to be hunted for too eagerly; insists
the Senate has now in fact bestowed it; and asks of Cato
the very thing Cato had asked of him --- to rejoice, if
what was preferred has come to pass, in its coming to
pass. The pointed touch is the closing observation that
Cato's friends were present at the drafting of the
decree, since (as everyone knows) such decrees are
drafted by the closest friends of the honoree. The
irritation is in the architecture, not in the prose;
the surface is courteous throughout. The letter closes
with the fear, already gathering in the homeward
letters, that the commonwealth Cicero is returning to
will not be in the state he wishes to find it.}
